\documentclass[11pt,a4paper]{article}

\usepackage[margin=1in]{geometry}
\usepackage{booktabs}
\usepackage{multirow}
\usepackage{hyperref}
\usepackage{xcolor}
\usepackage{amsmath}
\usepackage{array}
\usepackage{natbib}

\hypersetup{
  colorlinks=true,
  linkcolor=blue!60!black,
  citecolor=blue!60!black,
  urlcolor=blue!60!black
}

\title{%
  Content Over Format: What Drives Analysis-Mode in \\
  Small Language Models, and How Grounding Recovers Execution%
}

\author{Sophie D.\ Neilson \\
  \small Independent Researcher \\
  \small ORCID: \href{https://orcid.org/0009-0001-2430-1743}{0009-0001-2430-1743}
}

\date{September 2026}

\begin{document}
\maketitle

\begin{abstract}
A prior report observed that a small language model, when a task is prefixed with a
structured description of a recurring decision, stops acting and instead analyzes
the description. That study used one model and one scenario, and it did not separate
the description's content from its form. This follow-up adds the controls it lacked.
We prefix a task with either a three-axis descriptive ``glyph'' or an
information-matched imperative rule, across four decision classes, and we sweep the
descriptive prefix across four open-weight models on the original class. The shift
away from action generalizes across models, and it is driven by content rather than the
descriptive format: a matched imperative suppresses action as much as the glyph, and the
descriptive format shows no consistent advantage over a plain rule with the same content;
where the two differ, the glyph tends to suppress action slightly more, not less. Control prefixes (a content-scrambled glyph, a wrong-class glyph) show the
suppression mechanism is decision-class-dependent. We then vary two things
independently while holding the content fixed: form (descriptive versus directive)
and grounding (a domain-free statement versus a domain-specific one). Grounding is
the lever. A domain-specific aid restores execution where a domain-free one
suppresses it, in both forms. Whether a domain-specific \emph{description} suffices
as well as a domain-specific \emph{command} is under-determined here: the command
reaches ceiling in all four model-and-class cells, so the two separate only where
the description falls short of ceiling, which happens once, on the smaller model and
the harder task, and there the command does better. All runs use one local
raw-completion inference rig; reproducing them needs a 24\,GB GPU and no paid
interface. We state the scope and the measurement limits, including the cells where
a baseline already acts and no effect is measurable.
\end{abstract}

\section{Introduction}

A prior report~\citep{neilson2026register} documented a dissociation between
comprehension and execution in Mistral-7B-Instruct. When the model received a task
scenario prefixed with a \emph{glyph}, a three-axis structured description of a
recurring decision class, it analyzed the description rather than acting on the
task. Across three legs of eight runs it produced no task execution under the glyph,
while producing the target action in most runs with no glyph present. The report
named this the \emph{analysis-mode register shift} and scoped the observation to one
glyph, one scenario, and one model. It named format cueing as one of three candidate
explanations and disclaimed a mechanism.

That report named two open questions in its own limitations. The first is breadth:
whether the shift is specific to Mistral-7B. The second is the harder one. The glyph
carries information about the decision, and it also carries a particular form: named
axes, written descriptively rather than as instructions. The study compared a glyph
against no glyph, never against a different delivery of the same information, so it
could not say whether the descriptive \emph{format} drove the shift or whether the
\emph{content} did, by whatever route it arrived.

This follow-up runs the missing controls and then goes past them. We hold the
information constant and vary the delivery. For each decision class we write an
\emph{imperative}: a set of directives stating the same failure, correct action, and
applicability boundary as the glyph, in commanding rather than descriptive form,
audited for content parity. We run the glyph across four open-weight models. We add
two mechanism controls that damage the glyph in one way each: a \emph{scrambled}
glyph keeps the three-axis shape and the axis vocabulary but shuffles the word order
within each axis, destroying the propositional structure, and an \emph{off-target}
glyph is coherent but describes a different decision class. Finally we cross two
factors while holding content fixed: form (descriptive versus directive) and
grounding (a domain-free statement versus a domain-specific one).

The picture that results refines the prior observation rather than overturns it. The
register shift is real and general. It is carried by content, not chiefly by the
descriptive format, which adds only a small push toward suppression.
The content the glyph carries, present equally in a matched rule, is what moves
behavior; which part of the content matters depends on the decision class; and what
recovers execution is grounding the decision in the concrete domain.

\subsection{Terminology}\label{sec:terms}

These terms come from the research program this study belongs to, not from the
field's shared vocabulary. Each is defined here at first use.

\begin{itemize}
  \item \textbf{Glyph}: a three-axis structured description of a recurring decision
    class, written from inside the decision. Its axes are \emph{Marker} (what the
    failure looks like), \emph{Aim} (what correct navigation looks like), and
    \emph{Rest} (when the class does not apply). A glyph describes; it does not
    instruct. It is domain-free: it names abstract roles, not a specific tool or
    file.
  \item \textbf{Decision class}: a recurring situation in which an agent can act
    correctly or fail in a characteristic way, described apart from any one scenario.
  \item \textbf{Analysis-mode register shift}: the pattern in which a model, given a
    glyph, discusses the decision class rather than acting within it. It recognizes
    the obligation and does not carry it out.
  \item \textbf{Imperative}: a directive rewriting of a glyph's content. It states
    the same propositions as commands rather than as description, and stays
    domain-free. It is the content-and-form control.
  \item \textbf{Ground}: a domain-specific statement of the decision for the concrete
    scenario, in descriptive form (it says what a correct outcome is, without naming
    a command and without pasting the finished artifact).
  \item \textbf{Domain-directive}: the same domain-specific content as the ground, in
    directive form (it commands the action). Ground and domain-directive differ only
    in mood.
\end{itemize}

\subsection{Related Work}

Two recent papers occupy adjacent territory. Wang et al.~\citep{wang2026narrative}
study narrative priors: they run structurally identical games under different
storylines and find that the task narrative explains more behavioral variance than
an assigned persona does. Their work shows framing context shifts behavior. It does
not test a descriptive description of a decision class, and it does not separate
comprehension from execution. Liu et al.~\citep{liu2026behavioral} derive behavioral
mode axes, activation-space directions from contrastive behavioral traces, and use
them to measure and steer a model's behavioral style. Their notion of a behavioral
register is relevant to the analysis-mode phenomenon, though they steer through
activations rather than a prompt-level description.

The finding also sits next to a broader literature on prompt-format sensitivity and
on instruction-following in language models, which studies how the surface form of a
prompt changes behavior. We did not run a systematic survey of that literature, and
we make no universal novelty claim; a full positioning against it is future work.
The distinction this study turns on, description versus instruction, maps loosely
onto Ryle's \emph{knowing-that} versus \emph{knowing-how}~\citep{ryle1949}. The
theory behind the glyph draws on Polanyi's account of tacit
knowledge~\citep{polanyi1966} and Klein's recognition-primed decision
model~\citep{klein1998}, under which an agent that recognizes a decision class from
prior exposure navigates it without explicit rules. The present result bears on that
theory, and we return to it in the discussion.

\section{Materials and Method}

\subsection{Decision Classes}

We test four decision classes, each a domain-free description instantiated in one
task scenario.

\begin{itemize}
  \item \textbf{casg-direct}: an agent finishes a software release but leaves the
    companion \texttt{CHANGELOG.md} not updated. Correct action: add the release
    entry. This is the class from the prior report.
  \item \textbf{formal-step}: an agent must produce a compliance checklist by working
    it from the current specification, not by carrying a prior sign-off across.
    Working it surfaces a control the prior summary omits.
  \item \textbf{parent-state}: an agent must check a parent or milestone state before
    acting on a task.
  \item \textbf{conditional-gate}: a conditional gate that most models already handle
    unprompted in the tested scenario.
\end{itemize}

\subsection{Conditions}

The core study crosses two factors while holding the propositional content fixed,
giving a two-by-two plus a baseline:

\begin{itemize}
  \item \textbf{baseline}: the scenario alone.
  \item \textbf{glyph}: descriptive, domain-free (the three-axis glyph).
  \item \textbf{imperative}: directive, domain-free (a rule with the same content).
  \item \textbf{ground}: descriptive, domain-specific.
  \item \textbf{domain-directive}: directive, domain-specific.
\end{itemize}

The four aids carry the same core proposition at two grounding levels; a
content-parity audit records that the glyph and the imperative state the same facts,
that the ground and the domain-directive state the same facts and differ only in
mood, and that neither the ground nor the domain-directive hands over the finished
artifact. Two mechanism controls take the glyph's slot: a \emph{scrambled} glyph
(three-axis shape and axis vocabulary kept, word order shuffled within each axis, so the
propositional structure is destroyed) and an \emph{off-target}
glyph (a coherent glyph for a different class).

\subsection{The No-Tools Notice}\label{sec:notice}

The subject models are agentic: given a task, a strong model often emits tool calls
rather than inline prose. The probe runs in a single-turn harness that returns
nothing to a tool call, so a tool-calling reply stalls and produces no scoreable
action. That stall is a harness artifact rather than an analysis-mode response. We
append a short notice to the user turn, stating that no tools are available and to
answer inline. The notice is identical across all conditions, so it cannot bias the
contrasts, though it plausibly nudges every condition toward inline prose. Under the
notice the grids ran with no tool stalls.

\subsection{Models and Sampling}\label{sec:sampling}

The two-by-two ran on Qwen3.8-27B-Instruct (Q4\_K\_M) and Mistral-7B-Instruct-v0.3
(Q4\_K\_M), n=24 per cell. The four-class content grid and the mechanism controls ran
on the same two models at n=8 per cell. The cross-model sweep added phi-4-14B and
Qwen2.5-32B-Instruct on casg-direct at n=8.

All subjects are served by one local server (\texttt{llama-server}, llama.cpp) over a
raw completion path with a minimal per-model instruct wrapper, thinking disabled, and
no system prompt. One model runs at a time, verified from the server's properties
endpoint on each swap. Running every subject on this one local rig is a
reproducibility decision: the results reproduce on a single 24\,GB GPU with no paid
interface. A hosted model cannot run on this path, so no hosted model is a subject
here (Section~\ref{sec:scoring}).

The complete sampler chain, identical across conditions and models: temperature 0.8,
top\_k 0, top\_p 1.0, min\_p 0.05, all penalties disabled, maximum 1024 new tokens
(raised to 1536 for formal-step in the two-by-two, whose replies are longer; the
four-class grid and the mechanism controls used 1024), seeds 1 through the run count.
Penalties are disabled by design: a repetition penalty suppresses repeated structure,
and axis-by-axis analysis is repeated structure, so a live penalty would attenuate
the behavior under study. Qwen2.5-32B was served at a reduced 8K context to fit the
GPU, a KV-cache memory limit rather than a prompt-length one; its prompts run to about
1{,}200 tokens, well within the window. Each run records the served model, the build id,
the throughput, the sampler read back from the server, the image content digest, the
substrate, and the repository commit. The two-by-two ran on image
\texttt{sha256:67d26e73}; the mechanism controls on \texttt{sha256:e159c2b8}; the
four-class grid on \texttt{sha256:82182b74}; the sweep's new arms (phi-4,
Qwen2.5-32B) on \texttt{sha256:e159c2b8} and its Mistral and Qwen3.8 arms on
\texttt{sha256:82182b74}. Because the sweep pools cells from two images, one of which
followed a fix to a reply-truncation bug, we read the sweep as cells and direction,
not pooled counts; its replies are short, so the truncation fix does not bear on that
class.

\subsection{Scoring}\label{sec:scoring}

Responses are scored with a five-code rubric: \textbf{C}, recognition and the correct
action; \textbf{Ii}, recognition without action; \textbf{Ic}, recognition with a
defective action; \textbf{I}, no recognition; \textbf{N}, not scoreable. The
correct-action bar is identical across the glyph, imperative, ground, and
domain-directive conditions, so no scoring asymmetry can confound the contrasts. We
report an execution rate (an artifact produced: C or Ic) alongside the strict
correct-target rate (C), because the load-bearing distinction is whether the model
acts at all versus analyzes without acting. The per-class rubrics are in
Appendix~\ref{app:rubric}, and the strict correct-target table is
Table~\ref{tab:strict}.

Each cell was scored by two independent raters, both Claude (Anthropic) subagents,
blind to the condition and to the predictions, then adjudicated by the analyst.
Claude is a permitted judge but never a subject: it cannot run on the local
raw-completion rig that the reproducibility apparatus requires. The two-by-two (480
responses) reached a pooled pre-adjudication agreement of 477 of 480 (99.4\%); the
three disagreements, all in one model-and-class cell, were adjudicated against the
rubric. The predictions were written before the runs and kept out of everything the
subject and the rater saw. Two raters drawn from one model measure consistency, not
validity, because they share a training distribution. To test validity, an author
(a person) independently blind-scored a stratified 40-response subset, weighted to
the load-bearing casg-Mistral ground cell, blind to the condition and to the
automated codes. The author wrote the study's predictions, so this human anchor was
not blind to them; it was blind only to the condition and the automated codes. Human-to-Claude exact-code agreement was 39 of 40 (97.5\%), and 23
of 24 (96\%) on the load-bearing cell, where the human counted 12 executions to the
automated 13. The Claude scoring is therefore validated against a human anchor, and
the conclusions rest on gaps far larger than the one-run human-machine difference or
any plausible rater slip. Table~\ref{tab:agents} lists the raters.

\subsection{Agents and Models}\label{sec:agents}

\begin{table}[h]
\centering
\small
\caption{Agents and models. ``Not recorded'' means the study record does not hold the
  version. Run dates are UTC, matching the run-record timestamps; scoring and
  adjudication dates are the study's local recording dates and may differ by up to a
  calendar day.}
\label{tab:agents}
\begin{tabular}{@{}>{\raggedright\arraybackslash}p{3.2cm}>{\raggedright\arraybackslash}p{3.4cm}>{\raggedright\arraybackslash}p{4.8cm}>{\raggedright\arraybackslash}p{2.4cm}@{}}
\toprule
Role & Model & Version and runtime & Date \\
\midrule
Subject (two-by-two, grid, controls) & Qwen3.8-27B-Instruct; Mistral-7B-Instruct-v0.3 & Q4\_K\_M GGUF, llama.cpp raw completion, digest-pinned image & grid 2026-09-08; controls 2026-09-09; two-by-two 2026-09-15 \\
Subject (sweep) & phi-4-14B; Qwen2.5-32B-Instruct & Q4\_K\_M GGUF, llama.cpp; Qwen2.5-32B at 8K context & 2026-09-09 \\
Raters, two-by-two & Claude (Anthropic) subagents & Not recorded. Two blind raters per cell; adjudicated by the analyst. & 2026-09-14 \\
Raters, grid / sweep / controls & Not recorded & Blind subagents, model not named in those records. & 2026-09-08; 2026-09-09 \\
Human rater (validity anchor) & Author (a person) & Blind, stratified 40-response subset; 97.5\% exact-code agreement with the Claude raters. & 2026-09-14 \\
\bottomrule
\end{tabular}
\end{table}

\section{Results}

\subsection{The Register Shift Generalizes Across Models}\label{sec:sweep}

Table~\ref{tab:sweep} reports the correct-action count for casg-direct across four
models. Every model that acts unprompted at baseline moves away from action under a
glyph: Mistral 8 to 0, phi-4 8 to 0, Qwen3.8 8 to 4. On Mistral the suppressed runs
split evenly between recognition-without-action and no-recognition; on phi-4 they are
all recognition-without-action. The direction reproduces in three model families and
four sizes, so the analysis-mode shift is not specific to Mistral-7B.

\begin{table}[h]
\centering
\small
\caption{Correct-action count (out of 8) by model and condition, casg-direct.
  Non-C codes for the glyph and imperative are in parentheses.}
\label{tab:sweep}
\begin{tabular}{@{}lccc@{}}
\toprule
Model & Baseline & Glyph & Imperative \\
\midrule
Mistral-7B  & 8 & 0 (4 Ii, 4 I)  & 1 (7 Ii) \\
phi-4-14B   & 8 & 0 (8 Ii)       & 0 (8 Ii) \\
Qwen2.5-32B & 4 (4 Ii) & 2 (4 Ii, 2 Ic) & 6 (2 Ic) \\
Qwen3.8-27B & 8 & 4 (4 Ii)       & 4 (4 Ii) \\
\bottomrule
\end{tabular}
\end{table}

The same table gives the content answer for this class. The information-matched
imperative suppresses action as much as the glyph or more: identical at 0 for phi-4,
identical at 4 for Qwen3.8, near zero for Mistral, and, for Qwen2.5-32B, the plain
imperative reaches 6 where the glyph reaches 2. On this class the descriptive format
never beats the matched rule.

\subsection{Content Over Format in the Measuring Cells}\label{sec:grid}

The four-class content grid (n=8) extends the contrast to parent-state and
conditional-gate. The reading needs a calibration caveat: on the strong model the
baseline is at ceiling for casg-direct, formal-step, and conditional-gate, so those
cells cannot measure a guidance effect. Five of the eight class-and-model cells are at
ceiling; the contrast is measurable in three: parent-state on both models and
conditional-gate on Mistral. In them the glyph never beats the imperative. On
parent-state (Qwen) both lift correct action from a baseline of 1 of 8 to 8 of 8,
identically. On conditional-gate (Mistral) the imperative reaches 5 of 8 and the
glyph 2. On parent-state (Mistral) the imperative reaches 7 of 8 and the glyph 1,
because the model used the glyph's non-firing clause to reason itself out of the
check. The register shift measured as recognition without action is also not
format-specific: it appears under the glyph and the matched imperative at similar
rates (for example, casg on Qwen, 4 versus 4; formal-step on Qwen, 4 versus 5). The
four-model sweep of Section~\ref{sec:sweep} is casg-direct only, so the breadth is
four models on one class and four classes across three measuring cells, not a full
grid. Across the cells where the glyph and the imperative differ, the glyph more often
suppresses action slightly more than the matched imperative does, though not always:
the formal-step cells lean the other way. These cells span different grids, models, and
metrics, so pooling them is not clean, and a sign test on the lean is not robust to
which cells are counted. We read this as a weak directional tendency, not a measured
format effect: the descriptive format is not perfectly neutral, but it does not carry
the shift.

\subsection{The Suppression Mechanism Is Class-Dependent}\label{sec:mechanism}

Two control glyphs test what about the glyph does the work. A scrambled glyph keeps
the three-axis shape and the axis vocabulary but shuffles the word order within each
axis, destroying the propositional structure; an off-target glyph keeps coherence
and changes the class. On casg-direct both controls reproduce the shift (on Qwen the
scrambled glyph gives 5 recognition-without-action and the off-target 4, against the
real glyph's 4 and a baseline of 1), so on that class neither the propositional content
nor scenario recognition is necessary; the three-axis shape, with the class vocabulary
that survives the shuffle, carries it. That
reading does not generalize. On formal-step the scrambled glyph does not reproduce
the shift (the model acts, like baseline) while the off-target one does, so
comprehension is load-bearing there. On parent-state the off-target glyph fails to
cue the check while the scrambled one still cues it, so recognition is load-bearing
there. Across the informative classes, either comprehension or recognition carries
the effect, and not the same one. The honest statement is that the glyph's
suppression is class-dependent, and the control that fails names the ingredient the
class relies on. This walks back the single-class reading (that mere structure
suffices), which held only on casg-direct.

\subsection{Grounding Recovers Execution}\label{sec:ground}

The two-by-two (n=24) varies form and grounding with the content fixed.
Table~\ref{tab:ground} gives the execution rate. At n=24 a 95\% Wilson interval on a
proportion is roughly $\pm 0.18$ wide, so a gap smaller than about 5 of 24 is within
noise, and a gap of 10 or more of 24 clears it. That band is for a single proportion;
the uncertainty on a difference between two independent proportions is wider, by about
$\sqrt{2}$, so we give exact tests and difference intervals for the key contrasts. On
casg-Mistral, a Fisher exact test separates the domain-free glyph (0 of 24) from the
grounded description (13 of 24) at $p \approx 3\times10^{-5}$, and the grounded
imperative (5 of 24) from the domain-directive (24 of 24) at $p \approx 7\times10^{-9}$;
the glyph (0) and the domain-free imperative (5) differ at only $p \approx 0.05$, at the
edge of the band. A Newcombe 95\% interval on the imperative-to-domain-directive
execution difference is $+0.79$ $[+0.55, +0.91]$, and on the glyph-to-ground difference
$+0.54$ $[+0.31, +0.72]$. In every cell the domain-free aids
(glyph, imperative) suppress execution below baseline, and the domain command
(domain-directive) restores it to ceiling. The descriptive ground restores execution
relative to the domain-free aids in every cell (for example casg-Mistral, glyph 0 and
imperative 5 rise to ground 13), and it reaches ceiling on three of four cells. The
property that flips a prefixed decision-aid from suppressing action to permitting it
is grounding: a domain-free aid suppresses, a domain-specific aid recovers, in both
descriptive and directive form.

\begin{table}[h]
\centering
\small
\caption{Execution rate (an artifact produced: C or Ic, out of 24) in the
  form-by-grounding design.}
\label{tab:ground}
\begin{tabular}{@{}lcccc@{}}
\toprule
 & casg Mistral & casg Qwen & formal Mistral & formal Qwen \\
\midrule
baseline                         & 19 & 23 & 24 & 22 \\
glyph (descriptive, domain-free) & 0  & 15 & 14 & 16 \\
imperative (directive, domain-free) & 5 & 17 & 10 & 17 \\
ground (descriptive, domain-specific) & 13 & 24 & 24 & 24 \\
domain-directive (directive, domain-specific) & 24 & 24 & 24 & 24 \\
\bottomrule
\end{tabular}
\end{table}

Whether a domain-specific \emph{description} suffices as well as a domain-specific
\emph{command} is under-determined by this design. The domain-directive reaches
ceiling (24 of 24) in all four cells, so the two aids can be separated only where the
ground falls short of ceiling. That happens in one cell, casg on Mistral, and there
the command (24) beats the description (13) by a gap that clears the noise band. In
the other three cells both aids are at joint ceiling and cannot be distinguished. One
piece of evidence points the other way: on casg-Qwen the descriptive ground reaches
the strict correct-target bar in 24 of 24 runs against the command's 17
(Table~\ref{tab:strict}), so the description produced better-formatted output there.
Taken together, a grounded description clearly recovers execution, but the claim that
it does so \emph{as well as} a command is suggested, not established: the sole
discriminating cell favors the command.

\section{Discussion}

\subsection{The lever is grounding, not format or form}

The prior report named the glyph's descriptive format as one candidate cause of the
register shift. Under control that candidate does not carry the effect. A domain-free
imperative, the same content in directive form, induces the shift as strongly as the
glyph, so form is not the lever. Damaging the glyph shows the suppression runs
through different ingredients in different classes. And restoring execution turns on
one thing across every cell: whether the aid is grounded in the concrete domain. An
abstract decision-description, however phrased, pulls the model toward analyzing the
decision; grounding the same decision in the domain keeps the model acting on it. The
cleanest single line of evidence is the prose-versus-prose contrast, imperative to
domain-directive, which holds form constant and adds only grounding: execution rises
5 to 24, 17 to 24, 10 to 24, and 17 to 24 across the four cells.

\subsection{What this says about comprehension and compliance}

The recognition-primed account predicts that a model recognizing a decision class
navigates it without explicit rules. This study can say that a grounded description,
which commands nothing, recovers execution well above the domain-free aids, and
reaches ceiling in three of four cells. It cannot say that a description suffices as
well as a command, because the command is at ceiling everywhere and the two separate
in only one cell, where the command wins. So grounded comprehension is enough to
recover much of the execution the abstract aids suppressed, and whether it fully
matches a directive is open, with the one decisive comparison favoring the directive
and one strict-format comparison favoring the description. We do not resolve it here.

\subsection{Relation to the prior report}

This study extends the prior report rather than corrects it. The register shift the
report described is real and now shown across four models. What the report could not
see, for want of a content control, is that the effect was not carried by the
descriptive format. An earlier probe in the same program also reported that adding a
ground recovers execution, but that probe added the ground \emph{alongside} the
glyph and its ground pasted the finished artifact, so the recovery could have been
the model copying an answer. Here the ground stands alone and names the action
without writing it, and the recovery survives, so the ground effect is neither a
copy artifact nor dependent on the glyph. What the de-confounding changes is the
mechanism: grounding is the lever.

\subsection{Limitations}\label{sec:limitations}

\begin{enumerate}
  \item \textbf{Scope.} The two-by-two covers two classes and two models; the breadth
    results add two models on one class and two more classes at n=8, measurable in
    three cells. Each class has one scenario and one authored glyph. A different
    scenario or glyph could behave differently.
  \item \textbf{Small samples and ceilings.} Cells are 8 or 24 runs, and several
    baselines and the domain command are at or near ceiling, which is why the
    description-versus-command comparison separates in only one cell. We read
    direction and gaps that clear the noise band, not single integers.
  \item \textbf{Rater validity anchored on one human, not a second model family.} Both
    automated raters were Claude subagents, blind and adjudicated, with 99.4\% pooled
    agreement (consistency). Their validity was checked against a human anchor: an
    author blind-scored a stratified 40-response subset and agreed with the automated
    codes on 39 of 40 (97.5\%), and on 23 of 24 in the load-bearing cell. The human
    anchor is a single author rather than several independent people or a second model
    family, was the author of the predictions and so not blind to them (only to the
    condition and the automated codes), and it covers a subset, not the full grid; the
    conclusions rest on gaps far larger than the residual disagreement (for example 0 or
    5 rising to 24).
  \item \textbf{A residual in one contrast.} The glyph is a three-axis block and the
    other aids are prose, so the glyph-versus-ground contrast mixes grounding with
    block structure. The imperative-versus-domain-directive contrast (prose versus
    prose) is the cleaner grounding read and shows the same thing; the mechanism
    controls bound the structure question separately.
  \item \textbf{No mechanism below the behavioral level.} We report which ingredient
    is necessary per class and which delivery recovers execution. We do not claim to
    know why a prefixed description cues analysis inside the model.
  \item \textbf{No neutral-prefix control.} We did not run a neutral, non-glyph,
    length-matched prose prefix, so we cannot fully exclude that any long abstract
    prefix pulls a small model toward commentary. The scrambled and off-target controls
    keep the three-axis shape, so they do not settle this; a dedicated neutral-prefix
    control is planned.
  \item \textbf{Grounding is confounded with naming the target.} The grounded aids
    state what a correct outcome is for the scenario, so ``grounding'' here bundles
    domain-specificity with specifying the target. A grounded aid that names the domain
    but withholds the target outcome would separate the two, and is planned.
  \item \textbf{Setup, not agency.} Subjects run single-turn raw completion under a
    no-tools notice, not in a multi-turn tool-using loop. The results speak to small
    language models in this setup, not to deployed agents; whether a minimal tool-using
    loop changes the shift is untested and planned.
\end{enumerate}

\section{Conclusion}

A structured description of a decision class prefixed to a task suppresses execution
into analysis across four open-weight models. Under an information-matched control the
effect is driven by content rather than the descriptive format; a plain rule with the
same content suppresses as hard, and where the two differ the glyph suppresses slightly
more. Control prefixes show the suppression is class-dependent, resting
on bare structure, on comprehension, or on scenario recognition in different classes.
Execution is recovered by grounding the decision in the concrete domain, in both
descriptive and directive form. Whether a grounded description recovers execution as
fully as a command is left open by ceiling effects, with the one decisive cell
favoring the command. We report the finding at the scope of the tested glyphs,
scenarios, and models, and invite replication on the one local rig the study runs on.

\subsection{Data Availability}

The complete study materials (scenarios, glyphs, imperatives, grounds, scoring
rubrics, per-run scores, run records, and control materials) are available at
\url{https://github.com/sophdn/corpos-lab} under the
\texttt{studies/matched-content-experiment} directory.

\bibliographystyle{plainnat}

\begin{thebibliography}{9}

\bibitem[Klein(1998)]{klein1998}
Klein, G.~A. (1998).
\newblock \emph{Sources of Power: How People Make Decisions}.
\newblock MIT Press.

\bibitem[Liu et~al.(2026)]{liu2026behavioral}
Liu, H., Liu, R., Xu, H., Han, J., Fang, S., Yan, S., Deng, H., Wang, G., and
  Zou, N. (2026).
\newblock Your LLM, Your Style: Behavioral Mode Axes for LLM Behavioral
  Control.
\newblock \emph{arXiv preprint arXiv:2608.10703}.

\bibitem[Neilson(2026)]{neilson2026register}
Neilson, S.~D. (2026).
\newblock Structured Phenomenological Descriptions Induce Analysis-Mode Behavior
  in a Small Language Model.
\newblock \emph{Zenodo}. \href{https://doi.org/10.5281/zenodo.22542746}{doi:10.5281/zenodo.22542746}.

\bibitem[Polanyi(1966)]{polanyi1966}
Polanyi, M. (1966).
\newblock \emph{The Tacit Dimension}.
\newblock Doubleday.

\bibitem[Ryle(1949)]{ryle1949}
Ryle, G. (1949).
\newblock \emph{The Concept of Mind}.
\newblock Hutchinson.

\bibitem[Wang et~al.(2026)]{wang2026narrative}
Wang, Y., Lester, J., and Srivastava, S. (2026).
\newblock The Story Shapes the Agent: Narrative Priors in LLM Behavior.
\newblock \emph{arXiv preprint arXiv:2607.18566}.

\end{thebibliography}

\appendix

\section{Scoring Rubrics}\label{app:rubric}

The correct-action bar is condition-independent within a class. For casg-direct, C is
a \texttt{CHANGELOG.md} entry for v1.5.0 recording both changes under the appropriate
headings; an entry with the wrong file, a malformed version header, or one change
missing is Ic; recognition of the obligation with no entry is Ii; a conclusion that
the release is already complete is I. For formal-step, C is a worked v2 checklist that
includes the partner-pathway provenance control the prior summary omits; a worked
checklist that misses that control is Ic; naming the need without a worked checklist
is Ii; carrying the prior sign-off across as sufficient is I. For parent-state, C is a response
whose leading action confirms the parent milestone state (it consults the milestone
record to establish the current in-scope and scope-change state) before any
ticket-specific work; naming the check but starting ticket work first is Ii; consulting a
wrong parent artifact is Ic; diving into the task with no reference to the parent state is
I. For conditional-gate, C is a response that evaluates the pending-security-review
condition and withholds the merge on it; naming the condition but merging anyway is Ic,
and leaving the decision ambiguous is Ii; merging on green CI with no mention of the
condition is I; withholding by a uniform reflex without evaluating the specific condition
is also Ic. The full rubrics are committed with the study.

\section{Strict Correct-Target Rate}\label{app:strict}

Table~\ref{tab:strict} gives the strict correct-target rate (C only) for the
two-by-two, alongside the execution rate of Table~\ref{tab:ground}. On Mistral
casg-direct the baseline reaches the strict bar only once in 24 (it produces entries
in a wrong format), which is why the casg-Mistral recovery is read on the execution
rate; on casg-Qwen the descriptive ground reaches the strict bar more often than the
command (24 versus 17). This also reconciles an apparent gap with
Table~\ref{tab:sweep}, where the casg-Mistral baseline is 8 of 8: that count is from the
n=8 sweep on image \texttt{sha256:82182b74}, whose baseline produces correctly-formatted
entries, while the strict 1 of 24 here is from the two-by-two on image
\texttt{sha256:67d26e73}, whose baseline produces entries in a non-conforming format
(execution 19 of 24). The two grids ran on different images, so counts are read within
grid and by direction, not pooled across them.

\begin{table}[h]
\centering
\small
\caption{Strict correct-target rate (C only, out of 24).}
\label{tab:strict}
\begin{tabular}{@{}lcccc@{}}
\toprule
 & casg Mistral & casg Qwen & formal Mistral & formal Qwen \\
\midrule
baseline                         & 1  & 22 & 24 & 22 \\
glyph                            & 0  & 15 & 14 & 15 \\
imperative                       & 0  & 16 & 10 & 17 \\
ground                           & 12 & 24 & 24 & 24 \\
domain-directive                 & 24 & 17 & 24 & 24 \\
\bottomrule
\end{tabular}
\end{table}

\section{Materials}\label{app:materials}

The scenario, glyph, imperative, ground, and domain-directive for each class are
reproduced in the repository under \texttt{studies/matched-content-experiment}. The
glyph and imperative are content-parity audited against each other, and the ground
and domain-directive against each other, in the committed audit files.

\end{document}
